% ============================================================
% sec/05analisis_tecnico.tex — TRACK A
%
% OBLIGATORIO: análisis exploratorio profundo, diccionario de
% datos, calidad y estrategias de imputación justificadas,
% detección de valores atípicos, desbalance de clases, y
% prevención explícita de fuga de información.
%
% ESTADO: el contenido cuantitativo procede de la auditoría del
% conjunto de datos (Eje C). Este archivo establece la
% estructura y las justificaciones que dependen del Eje B.
% ============================================================
\section{Análisis Técnico del Track}
\label{sec:analisistecnico}

El \textit{track} aprobado corresponde al aprendizaje automático clásico
sobre datos tabulares. \pc{Fecha de aprobación.}

\subsection{Origen y composición del conjunto de datos}
\label{subsec:origen}

\pc{Descripción del conjunto: procedencia, periodo cubierto, número de
episodios, número de observaciones resultantes, y criterio de selección
aplicado por la fuente. Corresponde al Eje C.}

\noindent\textbf{Sesgo de selección.} \pc{Precisar el criterio de selección
de repeticiones aplicado por la fuente.} Saravanan y Guzdial observan que
el conjunto de estrategias dominantes refleja el comportamiento de la
mayoría de la base de jugadores, y orientan deliberadamente su agente a
aproximar un nivel promedio por esa razón
\cite{saravanan_metadiscovery_2024}. En la medida en que el conjunto
empleado aquí privilegie la franja superior de clasificación, la población
observada no corresponde a la población general.

\razon{Este sesgo debe declararse en el análisis técnico y recogerse en las
consideraciones éticas. Determina a qué población pueden extrapolarse los
resultados.}

\subsection{Análisis exploratorio}
\label{subsec:eda}

\pc{Número de registros y de características, distribución del objetivo,
cobertura temporal, distribución de la longitud de los episodios, y
estructura de agrupamiento. Corresponde a la auditoría del Eje C.}

\razon{Este apartado debe distinguir de forma inequívoca entre lo medido y
lo propuesto. Ninguna cifra procedente del diseño puede presentarse con la
forma de un resultado observado.}

\subsection{Diccionario de datos}
\label{subsec:diccionario}

\begin{table}[!t]
\centering
\caption{Diccionario de datos (extracto).}
\label{tab:diccionario}
\footnotesize
\begin{tabularx}{\columnwidth}{@{}P{1.6cm} P{1.1cm} P{1.5cm} Y@{}}
\toprule
\textbf{Variable} & \textbf{Tipo} & \textbf{Origen} &
\textbf{Significado y observaciones}\\
\midrule
\pc{var} & \pc{tipo} & \pc{origen} & \pc{descripción}\\
\bottomrule
\end{tabularx}
\end{table}

\pc{Diccionario completo o extracto representativo, con remisión al
repositorio para la versión íntegra.}

\noindent\textbf{Características derivadas por diferencia.} El diccionario
incorpora variables construidas como diferencia entre el valor
correspondiente al jugador de referencia y el de su oponente. Hodge
\textit{et al.} adoptan la misma construcción, calculando para cada métrica
el valor de cada bando y su diferencia \cite{hodge_winprediction_2021}.
Xenopoulos \textit{et al.} organizan igualmente su representación en
información global e información agregada por bando
\cite{xenopoulos_winprob_2022}.

\decidir{Hodge \textit{et al.} incorporan además, para cada característica,
el cambio respecto del instante inmediatamente anterior
\cite{hodge_winprediction_2021}. El diccionario actual no contempla
variables de esa naturaleza, que permitirían representar la variación
reciente del estado y no únicamente su valor. Su incorporación debe
evaluarse antes de materializar la tabla, verificando que la información
del instante anterior sea efectivamente anterior al instante de la
predicción y no constituya, por tanto, una fuga.}

\subsection{Calidad de los datos}
\label{subsec:calidad}

\pc{Verificaciones realizadas y su resultado: duplicados, consistencia de
esquema, rangos admisibles, consistencia lógica entre campos, y
cardinalidad de las variables categóricas. Corresponde a la auditoría del
Eje C.}

\subsection{Datos faltantes y estrategia de imputación}
\label{subsec:imputacion}

Se distinguen dos mecanismos de ausencia, con tratamientos distintos.

\noindent\textbf{Ausencia estructural.} Corresponde a campos ausentes
porque las reglas del juego determinan que no exista un valor en ese
estado. No constituyen datos faltantes en sentido estadístico y no se
imputan: se codifican de forma que la ausencia sea representable por el
modelo.

\noindent\textbf{Ausencia inesperada.} Corresponde a campos que deberían
tener valor y no lo tienen. \pc{Procedimiento de medición de su frecuencia
y criterio de decisión.}

\pc{Referencia arbitrada que sustente la distinción entre mecanismos de
ausencia. Corresponde al bundle de Daniel.}

\razon{La estrategia de imputación es una de las decisiones que el
enunciado exige justificar en el análisis técnico. La distinción entre
ambos mecanismos es correcta, pero requiere respaldo bibliográfico: sin él
constituye una decisión sin justificación técnica.}

Borisov \textit{et al.} señalan que los algoritmos modernos basados en
árboles manejan internamente los valores faltantes y los rangos extremos,
buscando aproximaciones y puntos de corte apropiados
\cite{borisov_tabular_2024}. Esta propiedad reduce la dependencia del
resultado respecto de la estrategia de imputación adoptada.

\subsection{Valores atípicos}
\label{subsec:atipicos}

\pc{Método de detección, umbral empleado, y tratamiento de las
observaciones detectadas.}

\razon{El criterio de exclusión debe declararse antes de examinar el
conjunto de prueba. Un valor extremo puede corresponder a una observación
válida y no a un error de registro; únicamente se excluye aquello
demostrablemente corrupto.}

\subsection{Desbalance de clases}
\label{subsec:desbalance}

\pc{Proporción medida sobre la representación tabular definitiva.}

\pc{Referencia arbitrada sobre el efecto de la corrección del desbalance
sobre las probabilidades predichas. Corresponde al bundle de Daniel.}

Silva Filho \textit{et al.} señalan que el desbalance de clases afecta las
probabilidades predichas y, en consecuencia, la calibración del modelo
\cite{silvafilho_calibration_2023}. Dado que la salida de este trabajo es
una probabilidad y no una etiqueta, cualquier procedimiento de corrección
del desbalance debe evaluarse por su efecto sobre la calibración y no
únicamente sobre las métricas de ordenamiento.

\razon{Este vínculo debe explicitarse. La decisión sobre el desbalance y la
decisión sobre la calibración no son independientes, y tratarlas por
separado ocultaría la interacción.}

\subsection{Prevención explícita de fuga de información}
\label{subsec:leakage}

Este apartado constituye el riesgo metodológico propio del \textit{track}
seleccionado.

\subsubsection{Información excluida del vector de características}

\pc{Enumeración de los campos excluidos y la razón de su exclusión en cada
caso.}

El criterio general es que el vector de características contenga
exclusivamente información disponible para un jugador en el instante de la
decisión. Xenopoulos \textit{et al.} enfrentan una restricción análoga
---aunque de origen distinto--- al limitar sus características a las
accesibles durante la transmisión en vivo, y reconocen que ello reduce el
conjunto disponible respecto de los trabajos que operan sobre repeticiones
completas \cite{xenopoulos_winprob_2022}. Hodge \textit{et al.} describen
la misma restricción, señalando que el flujo en vivo provee únicamente un
subconjunto de las características presentes en los archivos de repetición
\cite{hodge_winprediction_2021}.

\razon{Ambos precedentes son de origen técnico y el nuestro de origen
metodológico, pero el principio coincide: el modelo solo puede emplear
información que existiría en el momento real de la predicción. Citarlos
sitúa la restricción como práctica establecida.}

\subsubsection{Vectores de riesgo y su mitigación}

\begin{table}[!t]
\centering
\caption{Vectores de fuga identificados y su mitigación.}
\label{tab:leakage}
\footnotesize
\begin{tabularx}{\columnwidth}{@{}P{1.8cm} Y P{2.0cm}@{}}
\toprule
\textbf{Tipo} & \textbf{Descripción del riesgo} &
\textbf{Mitigación}\\
\midrule
Por episodio & Observaciones consecutivas de un mismo episodio difieren
mínimamente entre sí & Partición agrupada por episodio\\[2pt]
Temporal & Empleo de información posterior para predecir eventos
anteriores & Corte cronológico en el protocolo B\\[2pt]
Por agente & Identidad del agente disponible en entrenamiento y prueba &
\pc{mitigación}\\[2pt]
Por configuración de mazo & Memorización de qué configuraciones vencen &
Protocolo C\\[2pt]
Por información oculta & Inclusión de información no disponible para el
jugador & Exclusión explícita de los campos correspondientes\\[2pt]
Por metadatos & Identificadores u otros campos correlacionados con el
desenlace & Exclusión explícita\\[2pt]
Por preprocesamiento & Estimadores de transformación ajustados sobre la
tabla completa & Ajuste restringido al subconjunto de entrenamiento\\
\bottomrule
\end{tabularx}
\end{table}

\pc{Incorporar a la tabla la evidencia medida correspondiente a cada
vector. La cuantificación del riesgo de fuga por episodio bajo partición
aleatoria procede de la auditoría del Eje C.}

\razon{Figueiredo y Mendes cuantifican la contaminación entre subconjuntos
antes de entrenar modelo alguno \cite{figueiredo_leakage_2024}. Acompañar
cada vector con evidencia medida sobre este conjunto de datos, y no
únicamente con la descripción del riesgo, es lo que distingue una
prevención verificada de una precaución declarada.}

\subsubsection{Verificación automática}

\pc{Descripción de las pruebas automáticas que fallan ante la presencia de
un campo prohibido en el vector de características o ante la aparición de
un mismo episodio en dos subconjuntos.}

\razon{Una política de prevención sin una prueba que la verifique
constituye una intención y no una mitigación. La verificación automática es
lo que permite afirmar que la prevención se cumple y no solo que se
pretende.}

\decidir{Figueiredo y Mendes comparan explícitamente el desempeño obtenido
bajo distintas estrategias de partición, lo que permite cuantificar la
inflación atribuible a la fuga \cite{figueiredo_leakage_2024}. Un
procedimiento equivalente sobre este conjunto de datos ---ejecutar el
modelo de referencia bajo partición aleatoria y bajo partición agrupada, y
reportar la diferencia--- convertiría la prevención de fuga en un resultado
medido. Su incorporación debe decidirse en la Etapa 2, y en caso de
adoptarse debe declararse de forma inequívoca que la partición aleatoria se
emplea únicamente como demostración del sesgo y no para reportar
resultados.}